\section{Metric Spaces and Normed Spaces}

So far, we have been doing analysis on the real line \(\RR\), and occasionally on the complex plane \(\CC\). \(\RR\) has a structure that allows us to define what is meant by
\begin{itemize}
    \item Convergence of sequences, with unique limits when they exist;
    \item Cauchy sequences, which always converge (the completeness property of \(\RR\));
    \item Closed bounded intervals \(\ccintv{a}{b}\) which are special due to the Bolzano-Weierstrass theorem.
\end{itemize}

With this structure, we studied functions \(\functiondomain{f}{E \subseteq \RR}{\RR}\), where we study
\begin{itemize}
    \item The most basic notion, continuity of a function;
    \item If \(E = \ccintv{a}{b}\) and \(f\) is continuous, then \(f\) is bounded and attains its bounds (the extreme value theorem, from \textit{IA Analysis I});
    \item Furthermore, in the same scenario, \(f\) is uniformly continuous.
\end{itemize}

We would like to extend these structures to a more general setting of functions \(\functiondomain{f}{X}{\RR}\), or even functions \(\functiondomain{f}{X}{Y}\). The easiest way to impose a structure on \(X\) is define a notion of distance between points in \(X\), which is known as \term{metric}. More specifically, when \(X\) is already a (real) vector space, through the notion of ``norm'', we can define a metric on \(X\) that is compatible with the vector space structure.

\begin{definition}[Metric Space]
    Let \(X\) be any set. A \term{metric} (or \term{distance function}) on \(X\) is a function \(\functiondomain{d}{X \times X}{\RR}\) satisfying, for any \(x, y, z \in X\),
    \begin{enumerate}[label=(\roman*)]
        \item \(d\) is \term{non-negative}: \(d\pqty{x, y} \geq 0\), and \(d\pqty{x, y} = 0\) if and only if \(x = y\);
        \item \(d\) is \term{symmetric}: \(d\pqty{x, y} = d\pqty{y, x}\);
        \item \(d\) satisfies the \term{triangle inequality}: \(d\pqty{x, z} \leq d\pqty{x, y} + d\pqty{y, z}\).
    \end{enumerate}

    The pair \(\pqty{X, d}\) is called a \term{metric space}.
\end{definition}

\begin{definition}[Normed Vector Space]
    Let \(V\) be a vector space over \(\RR\). A \term{norm} on \(V\) is a function \(\functiondomain{\norm{\cdot}}{V}{\RR}\) satisfying, for any \(x, y \in V\) and \(\lambda \in \RR\),
    \begin{enumerate}[label=(\roman*)]
        \item \(\norm{\cdot}\) is \term{non-negative}: \(\norm{x} \geq 0\), and \(\norm{x} = 0\) if and only if \(x = 0\);
        \item \(\norm{\cdot}\) is \term{absolutely homogeneous} (compatibility with scalar multiplication): \(\norm{\lambda x} = \abs{\lambda} \norm{x}\);
        \item \(\norm{\cdot}\) satisfies the \term{triangle inequality}:
        \[
        \norm{x + y} \leq \norm{x} + \norm{y}.
        \]
    \end{enumerate}

    The pair \(\pqty{V, \norm{\cdot}}\) is called a \term{normed vector space} (or simply a \term{normed space}).
\end{definition}

\begin{proposition}[Normed Space induces a Metric Space]
    \label{prop:normed-space-induces-metric-space}%
    If \(\pqty{V, \norm{\cdot}}\) is a normed space, and if \(\functiondomain{d}{V \times V}{\RR}\) is defined by \(d\pqty{x, y} = \norm{x - y}\), then \(\pqty{V, d}\) is a metric space.
\end{proposition}

\begin{proof}
    We check that this induced metric is indeed a metric.
    
    \begin{enumerate}[label=(\roman*)]
        \item This follows from the non-negativity of the norm, since \(d\pqty{x, y} = \norm{x - y} = 0\) if and only if \(x - y = 0\), i.e., \(x = y\).
        \item This follows from the absolute homogeneity of the norm, since
        \[
            d\pqty{x, y} = \norm{x - y} = \norm{-(y - x)} = \abs{-1} \norm{y - x} = \norm{y - x} = d\pqty{y, x}.
        \]
        \item This follows from the triangle inequality for norms. Note
        \begin{align*}
            d\pqty{x, z} &= \norm{x - z} \\
            &= \norm{\pqty{x - y} + \pqty{y - z}} \\
            &\leq \norm{x - y} + \norm{y - z} \\
            &= d\pqty{x, y} + d\pqty{y, z}.
        \end{align*}
    \end{enumerate}
\end{proof}

\begin{example}[Finite-Dimensional Normed Spaces]
    The prototypical finite-dimensional vector space is \(\RR^n\).
    \begin{enumerate}
        \item We equip it with the Euclidean (or \(\ell_2\)) norm. We define the norm to be
        \[
            \norm[\text{E}]{x} = \norm[\ell_2]{x} = \norm[2]{x} = \pqty{\sum_{i = 1}^{n} \abs{x_i}^2}^{\flatfrac{1}{2}}.
        \]

        The only non-trivial property to check is the triangular inequality. Let \(x = \pqty{x_1, \dots, x_n}\) and \(y = \pqty{y_1, \dots, y_n}\). Then
        \begin{align*}
            \norm[2]{x + y}^2 &= \sum_{i = 1}^{n} \pqty{x_i + y_i}^2 \\
            &= \norm[2]{x}^2 + \norm[2]{y}^2 + 2 \sum_{i = 1}^{n} x_i y_i \\
            &\leq \norm[2]{x}^2 + \norm[2]{y}^2 + 2 \norm[2]{x} \norm[2]{y} && \text{Cauchy--Schwarz Inequality}\\
            &= \pqty{\norm[2]{x} + \norm[2]{y}}^2.
        \end{align*}

        Taking the square root finishes the proof.

        \item Another norm is the \(\ell_1\) norm, defined by
        \[
            \norm[\ell_1]{x} = \norm[1]{x} = \sum_{i = 1}^{n} \abs{x_i}.
        \]

        \item Another norm is the \(\ell_\infty\) norm, defined by
        \[
            \norm[\ell_\infty]{x} = \norm[\infty]{x} = \sup_{1 \leq i \leq n}\abs{x_i}.
        \]
    \end{enumerate}

    The proofs to that the final two norms satisfying the triangle inequality are left as exercises to the reader. In general, \(\pqty{\RR^k, \norm[p]{\cdot}}\) is a normed space for any \(p \geq 1\).
\end{example}

\begin{example}[Space of Real-Number Sequences]
    Denote \(\RR^\NN\) be the space of real-number sequences, i.e.
    \[
        \RR^\NN = \Bqty{\text{sequences }\pqty{x_k}_{k \in \NN}\text{ with }x_k \in \RR}.
    \]

    Under term-wise addition and scalar multiplication, this is a vector space.

    \begin{enumerate}
        \item Define the subspace
        \[
            \ell_1 = \set{\pqty{x_k} \in \RR^\NN}{\sum_{k = 1}^{\infty} \abs{x_k} < \infty}.
        \]

        This is a vector subspace of \(\RR^\NN\), which is still infinite dimensional, and we define the \(\ell_1\) norm on the \(\ell_1\) space by
        \[
            \norm[1]{x} = \sum_{k = 1}^{\infty} \abs{x_k}.
        \]
        \item Similarly, we may define
        \[
            \ell_2 = \set{\pqty{x_k} \in \RR^\NN}{\sum_{k = 1}^{\infty} \abs{x_k}^2 < \infty}.
        \]

        We define the \(\ell_2\) norm (the Euclidean norm) on the \(\ell_2\) space by
        \[
            \norm[2]{x} = \pqty{\sum_{k = 1}^{\infty} \abs{x_k}^2}^{\flatfrac{1}{2}}.
        \]
        \item Finally, we may define
        \[
            \ell_\infty = \set{\pqty{x_k} \in \RR^\NN}{\sup_{k \in \NN} \abs{x_k} < \infty}
        \]
        to be the space of bounded sequences. We define the \(\ell_\infty\) norm on the \(\ell_\infty\) space by
        \[
            \norm[\infty]{x} = \sup_{k \in \NN} \abs{x_k}.
        \]
    \end{enumerate}

    \(\ell_1, \ell_2, \ell_\infty\) are all infinite-dimensional normed spaces. The proofs are left to the reader.

    There is nothing special about the \(1, 2, \infty\) norms above. For any \(p \geq 1\), we may define the \(\ell_p\) norm to be
    \[
        \ell_p = \set{\pqty{x_k} \in \RR^\NN}{\sum_{k = 1}^{\infty} \abs{x_k}^p < \infty},
    \]
    and equipped with the \(\ell_p\) norm,
    \[
        \norm[p]{x} = \pqty{\sum_{k = 1}^{\infty} \abs{x_k}^p}^{\flatfrac{1}{p}},
    \]
    the space \(\pqty{\ell_p, \norm[p]{\cdot}}\) is a normed space. Checking the triangular inequality is non-trivial and is left to the reader.
\end{example}

\begin{example}[Space of Continuous Functions]
    Consider the vector space \(V = \contClass{0} \pqty{\ccintv{a}{b}}\) to be the space of real-valued continuous function on \(\ccintv{a}{b}\).

    \begin{enumerate}
        \item The \(L^1\) norm on \(V\) is defined by
        \[
            \norm[L^1\pqty{\ccintv{a}{b}}]{f} = \norm[L^1]{f} = \norm[1]{f} = \int_{a}^{b} \abs{f\pqty{x}} \dd{x}.
        \]

        \item The \(L^2\) norm on \(V\) is defined by
        \[
            \norm[L^2\pqty{\ccintv{a}{b}}]{f} = \norm[L^2]{f} = \norm[2]{f} = \pqty{\int_{a}^{b} \abs{f\pqty{x}}^2 \dd{x}}^{\flatfrac{1}{2}}.
        \]

        \item The \(L^\infty\) norm, also known as the \term{uniform norm} on \(V\) is defined by
        \[
            \norm[L^\infty\pqty{\ccintv{a}{b}}]{f} = \norm[L^\infty]{f} = \norm[\infty]{f} = \sup_{x \in \ccintv{a}{b}} \abs{f\pqty{x}}.
        \]
    \end{enumerate}

    These all define norms on \(V\). The triangular inequality for \(L^2\) follows from the integral version of the Cauchy--Schwarz inequality
    \[
        \int_{a}^{b} \abs{fg} \leq \pqty{\int_{a}^{b} \abs{f}^2}^{\flatfrac{1}{2}} \pqty{\int_{a}^{b} \abs{g}^2}^{\flatfrac{1}{2}}.
    \]

    To prove this, we use
    \[
        \phi\pqty{t} = \int_{a}^{b} \pqty{\abs{g} - t \abs{f}}^2
    \]
    and note that this quadratic in \(t\) has to be non-negative. The degenerate case where the integral of \(f\) is zero is trivial.
\end{example}

By \autoref{prop:normed-space-induces-metric-space}, all of the above examples are metric spaces with respect to the metric defined by the norm. For example, the \term{Euclidean metric} in \(\RR^n\) is induced by the \(\ell_2\) norm, as
\[
    d_{\text{E}} \pqty{x, y} = \norm[2]{x - y} = \pqty{\sum_{i = 1}^{n} \abs{x_i - y_i}^2}^{\flatfrac{1}{2}}.
\]

The \term{uniform metric} on \(\contClass{0}\pqty{\ccintv{a}{b}}\) is induced by the uniform \(L^\infty\) norm, as
\[
    d\pqty{f, g} = \norm[\infty]{f - g} = \sup_{x \in \ccintv{a}{b}} \abs{f\pqty{x} - g\pqty{x}}.
\]

\small

The ``natural'' space in which the integral norms \(L^1, L^2, \ldots\) should be defined is
\[
    \mathcal{R} \pqty{\ccintv{a}{b}}
\]
the space of integrable functions over \(\ccintv{a}{b}\). This is a metric space with respect to pointwise addition and scalar multiplication of functions.

There is an immediate concern here: for \(f \in \mathcal{R} \pqty{\ccintv{a}{b}}\), \(\norm[1]{f} = 0\) does not necessarily imply \(f = 0\) on \(\ccintv{a}{b}\) -- it only implies that \(f = 0\) almost everywhere on \(\ccintv{a}{b}\) by \autoref{thm:lebesgue-criterion-for-riemann-integrability}, i.e., \(f = 0\) except on a null set.

If \(\int \abs{f} = 0\) then \(f = 0\) at the points of continuities, and hence \(f = 0\) except on a null set (which is at most the set of discontinuities).

On the other hand, if \(f = 0\) almost everywhere on \(\ccintv{A}{b}\), then \(\int \abs{f} = 0\) by covering the null set with intervals of arbitrarily small total length, which is allowed by definition.

This allows us to define an equivalence relation \(f \sim g\) on \(\mathcal{R} \pqty{\ccintv{a}{b}}\) by \(f \sim g\) if and only if \(f = g\) almost everywhere on \(\ccintv{a}{b}\), i.e., the set of points at which they differ is a null set. So we define the quotient space \(\tilde{\mathcal{R}}\pqty{\ccintv{a}{b}} = \flatfrac{\mathcal{R} \pqty{\ccintv{a}{b}}}{\sim}\), which would become a genuine normed space, and indeed we can show that the induced operations on the equivalence classes are well-defined, and that the norm is well-defined on the equivalence classes since the integrals are equal.

\normalsize

\begin{definition}[Discrete Metric]
    For any set \(X\), the \term{discrete metric} on \(X\) is given by
    \[
        \function{d}{X \times X}{\RR}{\pqty{x, y}}{
            \begin{cases}
                0, & \qfor x = y, \\
                1, & \qfor x \neq y.
            \end{cases}
        }
    \]
\end{definition}

\begin{example}[New Metric from Given Metrics]
    Let \(\pqty{X, d}\) be a metric space. Then a new induced metric can be defined as
    \[
        \function{g}{X \times X}{\RR}{\pqty{x, y}}{\min\Bqty{1, d\pqty{x, y}}.}
    \]

    Another metric could be defined by
    \[
        \function{h}{X \times X}{\RR}{\pqty{x, y}}{\frac{d\pqty{x, y}}{1 + d\pqty{x, y}}.}
    \]

    Note that \(g\pqty{x, y} \leq 1\) and \(h\pqty{x, y} < 1\) for any \(x, y \in X\).
\end{example}

\begin{example}[French Railway Metric]
    The SNCF metric (also known as the French railway metric) is defined on \(\RR^2\). The idea is ``all trains go through Paris'' (which will be the origin here). Define
    \[
        d\pqty{x, y} = \begin{cases}
            \norm[2]{x - y}, & \qif x = ty \text{ for some }t \in \RR,\\
            \norm[2]{x} + \norm[2]{y}, & \quad\text{otherwise}.
        \end{cases}
    \]
\end{example}

\begin{definition}[Metric Subspaces]
    Let \(\pqty{X, d}\) be a metric space, and let \(Y\) be a subset of \(X\). Then the restriction \(\rightbar{d}_{Y \times Y}\) is a metric on \(Y\). This is known as the \term{subspace metric} or the \term{induced metric} on \(Y\).
\end{definition}

\subsection{Open and Closed Subsets}

\begin{definition}[Open and Closed Subsets]
    \label{def:open-and-closed-subsets}%
    Let \(\pqty{X, d}\) be a metric space.
    \begin{enumerate}[label=(\roman*)]
        \item Let \(a \in X\) be any point, with \(r > 0\). Then the \term{open ball with radius \(r\) and centre \(a\)} is the set
        \[
            \oball{r}{a} = \set{x \in X}{d\pqty{x, a} < r}.
        \]

        \item A subset \(U \subseteq X\) is \term{open} if for every point \(a \in U\), there exists \(r > 0\) such that \(\oball{r}{a} \subseteq U\).
        \item A subset \(F \subseteq X\) is \term{closed} if its complement \(X \setminus F\) is open.
    \end{enumerate}
\end{definition}